% Compiled with some variant of
% latexmk -pvc -shell-escape -halt-on-error
% Really just latexmk -pvc with my .latexmkrc file available at
% https://github.com/goedel-gang/dotfiles/blob/master/.latexmkrc

% Paper size, font size, equations align to the left
\documentclass[fleqn,a4paper,11pt]{article}
\title{``Guessing game'' assignment}
\author{Izaak van Dongen}

% Load all the configuration from a package, to keep content separate from
% markup.
\usepackage{mysty}

% The actual document content. Nowadays I would probably put this in a separate
% .tex file and \input that.
\begin{document}
    \maketitle\thispagestyle{empty} % no page number under title
    \tableofcontents
    \lstlistoflistings
    \listoffigures

    \section{Introduction}

    This is an assignment about guessing games and appropriate strategies. This
    is linked to searching algorithms, although in this case it's a sort of
    searching algorithm where you're searching for the item, rather than
    searching for the index of a known item.

    Both binary searches and linear searches can be adapted to that end, as both
    work by comparisons of the target and ``guesses''. These are analogous to
    asking the user questions - in the case of linear search, the question ``is
    your number equal to \ldots?'', and in the case of binary search, ``is your
    number greater than \ldots?'' Generalising from this idea of ``questions'',
    my general model/definition for a ``search algorithm'' is an algorithm that
    can ask only boolean questions of a user, and must use ideally as few of
    these as possible to guess a target number that the user knows.

    Therefore, essentially we need only to implement binary and linear search
    where the ``comparison function'' delegates to the user.  However,
    interpolation search, which has the complexity \(\BigO(\log(\log(n)))\) (on
    uniform data sets, which we can assume a user-generated target is), requires
    knowledge of the \emph{value} of the target, so it can be interpolated on
    the current upper and lower bounds. This means we will only achieve the
    \(\BigO(\log(n))\) complexity of binary search.

    Sadly, we also won't be able to implement any such search algorithm for the
    set of reals \(\mathbb{R}\), or any set with equivalent or higher cardinality
    (eg \(\mathbb{C} = \mathbb{R}^2 \implies \abs{\mathbb{C}} =
    \abs{\mathbb{R}}\)). This is because they are not a countably infinite set -
    there exists no bijection between them and the natural numbers. If we assume
    we could write a search algorithm to guess a real number, each of the
    algorithm's guesses could be labeled with a natural number, implying there
    does exist some enumeration of the reals. This contradicts lots of
    established knowledge (eg Cantor's diagonal argument) and so we know this
    won't be possible.

    Code for this assignment has been written in Pascal, targeted at and tested
    in fpc (version 3.0.0+dfsg-2), using mode \verb|{$MODE OBJFPC}|. This
    document was developed in \LaTeX.

    Where I've deemed it necessary, I've left comments, although often the code
    really is relatively simple, and the functions are sufficiently described by
    their names and the simplicity of the operations (``self-documented'').

    At some points, I use some algebraic notation. I use \(\land\), which means
    logical and, to indicate that multiple statements are true.

    Generally, I have provided pseudocode for my algorithms. My pseudocode is
    probably pretty close to structured english, but I feel like if I were to
    make it much more code-like it would more or less be Pascal/Python.

    \section{Boilerplate}

    To create a program that performs such a search to guess a number, some
    utility functions will be needed, that aren't really particularly
    algorithmic. I'm implementing these in a Pascal unit, named
    \verb|PUser.pas|. Having these in a separate unit means the code concerned
    with the actual algorithms can be more "purely" expressed. It also means
    that that code can become more modular itself, as another caller could
    feasibly implement their own boilerplate, and the abstracted algorithms will
    still work.

    \subsection{oracle}

    This is a function that can retrieve a boolean value from the user. This
    abstracts the ``boolean question'' part of the search algorithm away from
    the actual algorithmic logic.. It's named ``oracle'' as it acts as the
    comparison oracle the search functions must appeal to. The only parameter it
    takes is a message to display. The unix-inspired ruling it uses is that if
    an input starts with a ``y'', it is considered affirmative, and negative
    otherwise. The broad approach it uses is like this:

\begin{lstlisting}[caption=oracle pseudocode]
If the user's response is empty, the result is negative
If the user's response is not empty and begins with a lower or uppercase "y", the result is positive
Otherwise, the result is negative
\end{lstlisting}

    The catch for lower or uppercase is implemented by coercing the entire
    string to lowercase, using \verb|LowerCase|. It's implemented as:

\begin{lstlisting}[language=Pascal, caption=oracle implementation]
function oracle(msg: string): boolean;
var
    usr_input: string;
begin
    write(msg);
    readln(usr_input);
    if (length(usr_input) > 0) and (LowerCase(usr_input)[1] = 'y') then
        oracle := true
    else
        oracle := false;
end;
\end{lstlisting}

    NB: The lack of semicolon in line 8 is on purpose, as this is how an if-else
    statement works in Pascal (as it's considered one statement, the whole
    affair needs one semicolon).

    \subsection{read\_lohi}

    NB this function is currently unused in the source, but it was used for
    testing one of the earlier adaptations of the binary search algorithm. It
    also provides some interesting talking points about call signatures in
    Pascal.

    I wanted to provide some facility for setting the upper bound and lower
    bound of a binary search. As this function wants to kind of ``return'' two
    things, it has a nontrivial call signature, which is worth thinking about.
    One possible solution might be to define some kind of \verb|Bounds|
    container class that can hold both values.

    However, Pascal provides a way to output multiple values, and this is
    through an \verb|out| parameter. It could also be done using a \verb|var|
    parameter, although there's a slight difference. Both of these work, and in
    both cases the variable is passed by reference, but \verb|out| is slightly
    more specific - it's used when the input value of the variable is unneeded,
    which in this case is true, as I'm only concerned with using it as a channel
    for output. Once I've declared something as an \verb|out| parameter, I can
    assign to it and the caller can then use the new value (like the
    \verb|readln| function).

    I will be reading these from the command line arguments, as I'm developing
    it as a console application. As this is quite a simple application, I will
    just manually parse arguments with some if statements. The function will
    roughly do the following:

\begin{lstlisting}[caption=read\_lohi pseudocode]
Gather default values for the upper bound and lower bound
If a first command line argument is present
    Set the lower bound to this number
Otherwise, set it to the default
If a second command line argument is present
    Set the upper bound to this number
Otherwise, set it to the default
If either argument given was not an integer
    Set both to the defaults
\end{lstlisting}

    This is then implemented in Pascal using \verb|ParamStr| and
    \verb|ParamCount|. Interestingly, as it doesn't return anything in the
    conventional manner, this isn't a function but a procedure.

\begin{lstlisting}[language=Pascal, caption=read\_lohi implementation]
procedure read_lohi(lo_default, hi_default: integer;
                    out low_val, hi_val: integer);
begin
    try
        if ParamCount >= 1 then
            low_val := StrToInt(ParamStr(1))
        else
            low_val := lo_default;
        if ParamCount >= 2 then
            hi_val := StrToInt(ParamStr(2))
        else
            hi_val := hi_default;
    except
        on E: EConvertError do begin
            writeln('Conversion error occurred, reverting to defaults');
            low_val := lo_default;
            hi_val := hi_default;
        end;
    end;
end;
\end{lstlisting}

    \section{Linear Search}

    One approach to this is by a ``linear search''. This involved, basically, a
    kind of ``brute force'' approach - sequentially making guesses until one is
    correct. The single advantage of this algorithm is that it has no upper
    bound, and can even feasibly be made to work without a lower bound.

    \subsection{Positive integers}
    The most simple approach is to assume the number is some \(x\) such that
    \begin{equation}
        x \in \{0\} \cup \mathbb{N}
    \end{equation}
    ie the set of natural numbers including 0. Guesses can then be made
    sequentially like so:

\begin{lstlisting}[caption=Linear search on \(\mathbb{N}\) pseudocode]
Set the current ``guess'' to 0
While the guess is wrong
    Increment the guess by one
\end{lstlisting}

    In Pascal, this could be implemented like so:

\begin{lstlisting}[language=Pascal, caption=Linear search on \(\mathbb{N}\) implementation]
function linear_search: integer;
var
    i: integer = 0;
begin
    while not oracle('Is your number equal to ' + IntToStr(i) + '? ') do
        i := i + 1;
    linear_search := i;
end;
\end{lstlisting}

    \subsection{All integers}

    The flaw in the previous program is that if the user sneakily decides to
    think of a negative number, this program won't ever terminate. This can be
    solved by using some enumeration of the set of all integers \(\mathbb{Z}\).
    The simplest of these is the sequence 0, 1, -1, 2, -2, 3, -3, 4, -4, 5, -5,
    6, -6 \ldots This sequence doesn't duplicate 0, and has the property of
    ``spreading out'' from 0. A sequence kind of like this with a lot of
    overhead could be represented by simply generating the set of natural
    numbers \(\mathbb{N}\) and then taking each of these and its negation.
    However, if we consider it as more of a mathematical sequence, it
    \emph{can} be represented by a closed formula
    \begin{equation}
        U_n = (-1)^n \left\lfloor \frac{n}{2} \right\rfloor
    \end{equation}
    However, it's much better represented inductively with a condition:
    \begin{align*}
        U_1 &= 0 \\
        U_{n + 1} &=
        \begin{cases}
            -U_n + 1 & \text{if } U_n \leq 0 \\
            -U_n     & \text{if } U_n > 0
        \end{cases}
    \end{align*}
    This can also quite easily be implemented in code:

\begin{lstlisting}[language=Pascal,
                   caption=Linear search on \(\mathbb{Z}\) implementation]
function linear_search: integer;
var
    i: integer = 0;
begin
    while not oracle('Is your number equal to ' + IntToStr(i) + '? ') do
        if i <= 0 then
            i := -i + 1
        else
            i := -i;
    linear_search := i;
end;
\end{lstlisting}

    This does not need any double loops, duplicate code to explicitly negate
    each number, or extra boilerplate logic to prevent duplicating a 0.

    \subsection{Rational numbers}

    Interestingly, a linear search can also be generalised to work on the set of
    rational numbers \(\mathbb{Q}\). This is closely related to the fact that the
    cardinality of the set of rationals
    \begin{equation}
        \abs{\mathbb{Q}} = \abs{\mathbb{N}}
    \end{equation}
    This means there exists a bijection from \(\mathbb{N}\) to \(\mathbb{Q}\), and
    hence there must exist some enumeration of the elements of \(\mathbb{Q}\),
    which can be obtained by applying the bijection to all the naturals in
    order.

    We can easily come up with such an enumeration by considering the set of
    rationals as a grid with two integer axes representing numerator and
    denominator (in other words, the set \(\mathbb{Z}^2\)). We can then ``spiral
    outwards'' from the origin. This will be easiest to do in squares aligned
    with the axes, through which we move in the anticlockwise direction,
    although it could also be done in the diagonal direction.  The first part of
    the grid, in terms of co-ordinate pairs, will be like this:
    \begin{center}
        \begin{tabular}{ c c c c c }
            (-2, 2) & (-1, 2) & (0, 2) & (1, 2) & (2, 2) \\
            (-2, 1) & (-1, 1) & (0, 1) & (1, 1) & (2, 1) \\
            (-2, 0) & (-1, 0) & (0, 0) & (1, 0) & (2, 0) \\
            (-2, -1) & (-1, -1) & (0, -1) & (1, -1) & (2, -1) \\
            (-2, -2) & (-1, -2) & (0, -2) & (1, -2) & (2, -2) \\
        \end{tabular}
    \end{center}
    Where the first couple of terms of the enumeration are:
    \begin{equation}
        \begin{split}
        &(0, 0), \\
        &(1, 1), (0, 1), (-1, 1), (-1, 0), (-1, -1), (0, -1), (1, -1), (1, 0), \\
        &(2, 1), (2, 2), (1, 2), (0, 2), (-1, 2), (-2, 2), (-2, 1), (-2, 0), \ldots \\
        &(3, 1), (3, 2), \ldots
        \end{split}
    \end{equation}
    Note that the first term is \(\frac{0}{0}\). This technically is not a
    rational number as its value is undefined, but we can deal with that later
    - first we can establish the enumeration, and then ``filter'' it for
    unwanted pairings. Note also that this enumeration will not preserve any
    conception of size of fractions - \(\frac{1}{2}\) will occur long before
    \(\frac{1}{100}\).

    We know that this path will reach all of the points on this grid, hence this
    enumeration will eventually reach the target - and this is where a linear
    search is useful. First, we will need to consider how we can implement this
    enumeration. We can consider the desired behaviour on each side of a square,
    and the conditions describing each such side.

    Let us consider two sequences \(X\) and \(Y\), where we say that
    \(X_i / Y_i\) is the \(i\)th element of our enumeration.

    We know the first term is \(\frac{0}{0}\). Hence, we have \(X_1 = 0 \land Y_1
    = 0\)

    We also say that we jump to the next layer of square when we ``hit'' the
    positive x-axis. To start the next square, we move to the right one and up
    one, as if we just move to the right, we stay on the axis and move right
    again. This means
    \(X_i \geq 0 \land Y_i = 0 \implies X_{i + 1} = X_i + 1 \land Y_{i + 1} =
    1\)

    If we are on the right side of the square, we know that \(X_i\) is positive.
    We also know that \(-X_i \leq Y_i < X_i\). When this is true, we want to move
    ``up'', so increment \(Y\). We also know that \(Y_i = 0\) is considered a
    special case. Hence, we say that
    \begin{equation}
      X_i > 0 \land
      -X_i \leq Y_i < X_i \land
      Y_i \neq 0
    \implies
      X_{i + 1} = X_i \land
      Y_{i + 1} = Y_i + 1
    \end{equation}
    We can produce similar definitions for the other three sides:
    \begin{alignat*}{7}
    &Y_i &&> 0 &&\land -Y_i &&< X_i &&\leq Y_i &&\implies X_{i + 1} = X_i - 1 &&\land Y_{i + 1} = Y_i \\
    &X_i &&< 0 &&\land X_i &&< Y_i &&\leq -X_i &&\implies X_{i + 1} = X_i &&\land Y_{i + 1} = Y_i - 1 \\
    &Y_i &&< 0 &&\land Y_i &&\leq X_i &&< -Y_i &&\implies X_{i + 1} = X_i + 1 &&\land Y_{i + 1} = Y_i
    \end{alignat*}
    Now that we have a way to produce such an enumeration, we can consider a
    filtering. The first condition is that we skip any number with a
    denominator of 0, as this is undefined. The second condition is that we
    should only consider fractions in their simplest form - ie the numerator
    and denominator are coprime. This means we don't duplicate any rationals
    and so obtain a strict bijection between the naturals and the rationals.
    Knowing that
    \begin{equation}
        \text{A and B are coprime} \Leftrightarrow gcd(A, B) = 1
    \end{equation}
    We can simply calculate the gcd of \(X\) and \(Y\) using Euclid's algorithm,
    implemented like so in Pascal:

\begin{lstlisting}[language=Pascal, caption=Euclid's algorithm in Pascal]
function gcd(a, b: integer): integer;
var
  temp: integer;
begin
  while b <> 0 do begin
    temp := b;
    b := a mod b;
    a := temp;
  end;
  gcd := a;
end;
\end{lstlisting}

    This also works how we want for negative numbers - printing a couple of
    examples:

\begin{lstlisting}[caption=GCD behaviour for negative numbers]
gcd 1 2 is 1
gcd 1 -2 is 1
gcd -1 2 is -1
gcd -1 -2 is -1
\end{lstlisting}

    We can see that only one of the pairs \(1, -2\) and \(-1, 2\) are considered
    coprime, which is correct as \(\frac{1}{-2} = \frac{-1}{2} = -\frac{1}{2}\).
    We can also see that -1 and -2 are not considered coprime, which is again
    correct as it is equivalent to \(\frac{1}{2}\).

    We can now put all of this together to obtain a fully featured linear search
    algorithm on the set of rational numbers:

\begin{lstlisting}[language=Pascal, caption=Linear search on \(\mathbb{Q}\) implementation in Pascal]
procedure linear_search(out o_x: integer;
                        out o_y: integer);
var
    x: integer = 0;
    y: integer = 0;
begin
    while (y = 0) or (gcd(x, y) <> 1) or
      (not oracle('Is your number equal to ' + IntToStr(x) + '/' + IntToStr(y) + '? ')) do
        if (x >= 0) and (y = 0) then begin
            x := x + 1;
            y := y + 1;
        end else if (x > 0) and (-x <= y) and (y < x) then
            y := y + 1
        else if (y > 0) and (-y < x) and (x <= y) then
            x := x - 1
        else if (x < 0) and (x < y) and (y <= -x) then
            y := y - 1
        else
            x := x + 1;
    o_x := x;
    o_y := y;
end;
\end{lstlisting}

    This again uses the \verb|out| parameter trick to return multiple values. It
    is also a little terser than the fully qualified mathematical conditions
    given earlier, as some conditions are mutually exclusive/eliminated.
    Interestingly, the body of the while loop is technically a single statement,
    so no begin \ldots end block is needed. The condition of the loop is that it
    keeps going if the number is one to be skipped, or the user does not
    indicate that the guess is correct.

    Another thing this code makes use of is short-circuit boolean evaluation.
    All of the ``filter'' checks have been compounded into one boolean
    expression, using the or operator. Even though there is a call to oracle in
    this statement, the call will only actually be evaluated if neither of the
    filter checks flag up - if it doesn't pass the filter, it will short-circuit
    so the call to oracle will be skipped.

    A parting consideration is that the linear search is very ``secure'', in a
    sense - a malicious user will not be able to confuse the linear search in
    any way, they will only be able to jeopardise their own experience - any
    sequence of answers for a linear search will seem sensible to it, even if
    they are intended maliciously. Really, a user's only power is to lie about
    what their number was, which a simple command line program won't
    realistically be able to defend against.

    \section{Binary search}

    At last! The binary search algorithm is significantly faster than the
    linear search, in the average case (being \(\BigO(\log(n))\)).

    \subsection{Basic binary}

    In its most basic form, the binary search might work something like this in
    pseudocode:

\begin{lstlisting}[caption=Basic binary search pseudocode]
Establish a lower bound and a strict upper bound
While the difference between the bounds is greater than one
    Guess a number in the middle of these bounds
    If the guess was too high, move the upper bound down to the guess
    If the guess was too low, move the lower bound up to the guess
\end{lstlisting}

    Note that this algorithm will rely on some "magic" upper and lower bound,
    for now. I've used some slightly careful consideration of range arithmetic
    and definitions - by having a strict upper bound, when the length of the
    range is one, the only remaining candidate in the range is the lower bound.
    This is one easy way to prevent off-by-one errors and too much
    case-checking here. It can be implemented in Pascal quite trivially, like
    so:

\begin{lstlisting}[language=Pascal, caption=Basic binary search in Pascal, label={lst:binarysearch}]
function binary_search(lower, upper: integer): integer;
var
    mid: integer;
    is_ge: boolean;
begin
    while upper - lower > 1 do begin
        mid := (upper + lower) div 2;
        is_ge := oracle('Is your number greater than or equal to ' + IntToStr(mid) + '? ');
        if is_ge then
            lower := mid
        else
            upper := mid;
    end;
    binary_search := lower;
end;
\end{lstlisting}

    This function takes two parameters, the bounds, and returns the guessed
    number. It is also relatively immune to malicious users - again, every
    sequence of answers leads to a correct guess, assuming all the answers were
    correct.

    \subsection{Bound-finding binary search}

    Our binary search could be made significantly less na\"ive by not putting a
    bounds restriction on the user. We can actually determine a bound quite
    quickly, assuming nothing other than that the the user's number will be an
    integer.

    Continuing the theme of binary and powers of two, we can simply check the
    ``bound'' between each of the adjacent powers of 2, ie first 1-2, then 2-4,
    then 4-8 (note that these are the good old fashioned strict upper bound
    ranges, so there is no overlap). This will converge to some number \(n\) in
    logarithmic time, again. However, this does not account for negative
    numbers - what we can easily do is simply consider the \emph{absolute}
    value of the user's number \(\abs{n}\). We need then only ask if it is
    negative once, at the end. If it is negative, we will need to ``flip'' the
    range, bearing in mind to increase the new lower bound and decrease the new
    upper bound, to make them non-strict and strict, respectively. This leaves
    the single edge case where the user's number is 0, as this is not an
    integer power of 2. We can explicitly check for this to get around it (not
    very elegant, but simple enough).

    In pseudocode, this might be implemented as follows:

\begin{lstlisting}[caption=Bound-finding pseudocode]
If the user's number is 0, return the range 0-1
For each pair a and b of powers of 2
    If a <= |n| < b
        Remember this range, and stop guessing ranges
If the user's number is positive, return this range
Otherwise, return the ``flipped'' range
\end{lstlisting}

    This then comes to the following implementation in Pascal:

\begin{lstlisting}[language=Pascal, caption=Bound-finding subroutine in Pascal]
procedure binary_find_bounds(out o_lower: integer;
                             out o_upper: integer);
var
    lower: integer = 1;
    upper: integer = 2;
begin
    writeln('This is the stage where I determine some bounds on your number. I'
           , ' will be asking questions about the *magnitude* of your number,'
           , ' so watch out.');
    if oracle('Is your number 0? ') then begin
        o_lower := 0;
        o_upper := 1;
    end else begin
        while not oracle('Is your number some x such that ' +
          IntToStr(lower) + ' <= |x| < ' + IntToStr(upper) + '? ') do begin
            lower := lower * 2;
            upper := upper * 2;
        end;
        if oracle('Is your number greater than or equal to 0? ') then begin
            o_lower := lower;
            o_upper := upper;
        end else begin
            o_lower := - upper + 1;
            o_upper := - lower - 1;
        end;
    end;
end;
\end{lstlisting}

    Using this, we can successfully perform a binary search on the set of all
    integers \(\mathbb{Z}\). In fact, we only need to first call this subroutine
    to determine the bounds, and then let the previous function from listing
    \ref{lst:binarysearch} do the rest.

    \subsection{Rational numbers}

    Interestingly, theoretically we already have enough ``material'' to put
    together a binary search on the set of all rationals. We can use the
    previously defined enumeration, and then simply ask the user to think of
    what the index of their number is in that enumeration, and then perform a
    binary search as normal to find that integer. Here, the concept of ``greater
    than'' has been entirely abstracted away from the comparison of actual
    rationals.

    However, this makes some rather theoretical demands of the user, and is
    better left as a perverse thought experiment. There are far better ways to
    approach this. Namely, there exists a binary search tree of the set of
    rationals - the ``Stern-Brocot'' tree, which uses fractional mediants:

    \begin{figure}[h]
        \includegraphics[width=0.5\textwidth]{sternbrocot}
        \centering
        \caption{A Stern-Brocot tree}
    \end{figure}

    This also addresses the concern about comparison - this tree does preserve a
    proper ordering of the rationals, due to the fact that the mediant of two
    fractions lies strictly in between those fractions.

    Interestingly, this tree has kind of similar behaviour to what we
    encountered in the bound-finding function - it only terminates on positive
    nonzero numbers. We can get around this in the same way - first check for 0,
    ask questions about the absolute value of the user's number and then ask if
    it is negative:

\begin{lstlisting}[caption=Binary search on \(\mathbb{Q}\) in pseudocode]
Ask if the user's number is 0
Set the lower and upper bound to 0/1 and 1/0
While you haven't guessed the user's number
    Generate the mediant of the bounds
    If this is the number, stop
    If this is too high, move the upper bound down to the mediant
    Otherwise, move the lower bound up to the mediant
\end{lstlisting}

    This can be implemented like so in Pascal:

\begin{lstlisting}[language=Pascal, caption=Binary search on \(\mathbb{Q}\) in Pascal]
procedure binary_search(out o_x: integer;
                        out o_y: integer);
var
    mid_x, mid_y: integer;
    lo_x: integer = 0;
    lo_y: integer = 1;
    hi_x: integer = 1;
    hi_y: integer = 0;
    is_ge: boolean;
begin
    if oracle('Is your number equal to 0? ') then begin
        o_x := 0;
        o_y := 1;
    end else begin
        repeat
            mid_x := lo_x + hi_x;
            mid_y := lo_y + hi_y;
            is_ge := oracle('is your number q such that '
                          + IntToStr(mid_x) + '/' + IntToStr(mid_y)
                          + ' <= |q|? ');
            if is_ge then begin
                lo_x := mid_x;
                lo_y := mid_y;
            end else begin
                hi_x := mid_x;
                hi_y := mid_y;
            end;
        until is_ge and oracle('is your number q such that |q| = '
                          + IntToStr(mid_x) + '/' + IntToStr(mid_y) + '? ');
        o_x := mid_x;
        o_y := mid_y;
        if oracle('Is your number less than 0? ') then
            o_x := -mid_x;
    end;
end;
\end{lstlisting}

    As it is traversing an infinite tree, it will never reach a point where the
    bounds are so close that it can stop. Therefore, as it only has boolean
    questions it must sometimes ask if the user's number is exactly equal to a
    guess, to produce a terminating condition.

    Note that although it uses a binary search tree, it has diverged a little
    from the classic idea of a binary search in the context of a guessing game,
    as the way it ``splits'' choice is not by mean but by mediant.
    Interestingly, if we tried this approach but using the mean of the upper and
    lower bounds, it would never terminate for the vast majority of fractions -
    this would be because it could only ``reach'' fractions with a denominator
    which is a power of 2. Interestingly, this is also exactly the set of
    rationals representable by a float (although a finite float can of course
    only represent a subset of these).

    It's also not entirely user-friendly, in that if a user decides on a crazy
    fraction like \(\frac{51}{97}\), they have to do a lot of mental maths to
    answer the questions. However, really there isn't any way around this if you
    want to stay true to the binary search idea of eliminating many possible
    guesses simultaneously.

    \section{Interface}

    Having implemented all of the algorithms, it was now time to write an
    interface. First, I wrote a number of procedures to go with each algorithm,
    prefixed with ``perform\_''. They mainly just printed some information and
    formatted the results to the user. For example, to go with the binary search
    on \(\mathbb{Z}\):

\begin{lstlisting}[language=Pascal, caption=Binary search on \(\mathbb{Z}\) wrapper procedure]
procedure perform_bsearch;
var
    guess: integer;
    lower, upper: integer;
begin
    writeln('welcome to the binary search game! Think of an integer.');
    binary_find_bounds(lower, upper);
    guess := binary_search(lower, upper);
    writeln('your number was ', guess);
end;
\end{lstlisting}

    These procedures are all stored in Pascal \emph{units}, which provide an
    interface to the ``perform'' procedures. Using such procedures, I then wrote
    the following short program:

\begin{lstlisting}[language=Pascal, caption=Main interface loop]
while true do begin
    if oracle('Would you like to perform an integer binary search? ') then
        perform_bsearch
    else if oracle('Would you like to perform a rational binary search? ') then
        perform_fracbsearch
    else if oracle('Would you like to perform an integer linear search? ') then
        perform_lsearch
    else if oracle('Would you like to perform a rational linear search? ') then
        perform_fraclsearch;
end;
\end{lstlisting}

    Note that there is no explicit way to exit - this is because I favour the
    tried and tested command line idiom of the keyboard interrupt. When the user
    wishes to exit, they can simply hit Ctrl-c and get on with their day.

    \section{``Testing''}

    Here is an example of a session you might enter with the program:

\begin{lstlisting}[caption=Typical session]
To answer affirmatively, you should write anything starting with a "y". Anything else will be considered negative. Lower and upper bounds can be provided in argv
Would you like to perform an integer binary search? y
welcome to the binary search game! Think of an integer.
This is the stage where I determine some bounds on your number. I will be asking questions about the *magnitude* of your number, so watch out.
Is your number 0? n
Is your number some x such that 1 <= |x| < 2? n
Is your number some x such that 2 <= |x| < 4? n
Is your number some x such that 4 <= |x| < 8? n
Is your number some x such that 8 <= |x| < 16? y
Is your number greater than or equal to 0? y
Is your number greater than or equal to 12? n
Is your number greater than or equal to 10? y
Is your number greater than or equal to 11? n
your number was 10
Would you like to perform an integer binary search? n
Would you like to perform a rational binary search? y
Welcome to the binary search game! Think of a rational number.
Is your number equal to 0? n
is your number q such that 1/1 <= |q|? y
is your number q such that |q| = 1/1? n
is your number q such that 2/1 <= |q|? n
is your number q such that 3/2 <= |q|? n
is your number q such that 4/3 <= |q|? y
is your number q such that |q| = 4/3? n
is your number q such that 7/5 <= |q|? y
is your number q such that |q| = 7/5? y
Is your number less than 0? y
your number was -7/5
Would you like to perform an integer binary search? n
Would you like to perform a rational binary search? n
Would you like to perform an integer linear search? y
Think of an integer
Is your number equal to 0?
Is your number equal to 1?
Is your number equal to -1?
Is your number equal to 2?
Is your number equal to -2?
Is your number equal to 3?
Is your number equal to -3? y
Your number was -3
Would you like to perform an integer binary search?
Would you like to perform a rational binary search?
Would you like to perform an integer linear search?
Would you like to perform a rational linear search? y
Think of a rational number
Is your number equal to 1/1?
Is your number equal to 0/1?
Is your number equal to -1/1?
Is your number equal to 2/1?
Is your number equal to 1/2?
Is your number equal to -2/1?
Is your number equal to 1/-2? y
Your number was 1/-2
Would you like to perform an integer binary search? ^C
\end{lstlisting}

    \section{Source}
    The full project in its directory structure, including this document (as a
    full-colour PDF and \TeX{} file), can be found at
    \url{https://github.com/goedel-gang/assignment_guessing}.

\end{document}
